\documentclass[11pt,a4paper]{article}

% ============================================================
%  QIMMA -- COURS : Fonction logarithme népérien
%  2ème Bac Sciences Mathématiques (A et B)
%  Standard exhaustif : définitions formelles + démonstrations
%  Basé sur templates/qimma-template.tex (charte Qimma)
% ============================================================

\usepackage[french]{babel}
\usepackage[utf8]{inputenc}
\usepackage[T1]{fontenc}
\usepackage{lmodern}
\usepackage[a4paper,margin=2cm,top=2.3cm,bottom=2.6cm]{geometry}
\usepackage{amsmath,amssymb}
\usepackage{xcolor}
\usepackage{graphicx}
\usepackage{tikz}
\usetikzlibrary{shadows,calc}
\usepackage[most]{tcolorbox}
\usepackage{titlesec}
\usepackage{fancyhdr}
\usepackage{enumitem}
\usepackage{pifont}
\usepackage{array}
\usepackage{colortbl}
\usepackage{hyperref}

% ------------------- CHARTE QIMMA -------------------
\definecolor{qteal}{HTML}{0d9488}
\definecolor{qtealdark}{HTML}{134e4a}
\definecolor{qamber}{HTML}{fbbf24}
\definecolor{qambertext}{HTML}{92400e}
\definecolor{ink}{HTML}{1f2937}
\definecolor{inkgray}{HTML}{6b7280}
\definecolor{tealpale}{HTML}{e6f5f3}
\colorlet{colDef}{qteal}
\definecolor{colProp}{HTML}{0f766e}
\definecolor{colEx}{HTML}{b45309}
\definecolor{colRem}{HTML}{9d174d}
\color{ink}

\graphicspath{{../../../../assets/}}
\newcommand{\logofile}{../../../../assets/qimma-logo.pdf}

% ------------------- METADONNEES -------------------
\newcommand{\matiere}{Mathématiques}
\newcommand{\niveau}{2\up{ème} Bac -- Sciences Mathématiques}
\newcommand{\chapitretitre}{Fonction logarithme népérien}

% ------------------- EN-TETE / PIED DE PAGE -------------------
\pagestyle{fancy}
\fancyhf{}
\renewcommand{\headrulewidth}{0pt}
\renewcommand{\footrulewidth}{0.5pt}
\fancyfoot[L]{\raisebox{-0.15cm}{\includegraphics[height=0.35cm]{\logofile}}~\small\sffamily\color{inkgray}Qimma}
\fancyfoot[C]{\small\sffamily\color{inkgray}\matiere\ -- \niveau}
\fancyfoot[R]{\small\sffamily\color{qtealdark}\thepage}

% ------------------- TITRES DE SECTION -------------------
\titleformat{\section}
  {\normalfont\sffamily\Large\bfseries\color{qtealdark}}
  {\tikz[baseline=-3.2pt]{\node[circle,fill=qteal,text=white,inner sep=0pt,minimum size=8mm]{\sffamily\bfseries\thesection};}}
  {10pt}{}
  [\vspace{-2mm}{\color{qamber}\titlerule[1.6pt]}\vspace{2mm}]
\titlespacing*{\section}{0pt}{16pt}{10pt}

\titleformat{\subsection}
  {\normalfont\sffamily\large\bfseries\color{qtealdark}}
  {\color{qteal}\ding{212}}{6pt}{}
\titlespacing*{\subsection}{0pt}{12pt}{6pt}

% ------------------- BOITES PEDAGOGIQUES -------------------
\newtcolorbox{fiche}[3][]{
  enhanced, breakable,
  colback=white, colframe=#2,
  boxrule=1pt, arc=2mm,
  top=7mm, left=4mm, right=4mm, bottom=3mm,
  drop shadow={black!25, shadow xshift=1mm, shadow yshift=-1mm},
  attach boxed title to top left={xshift=4mm, yshift=-3mm, yshifttext=-1mm},
  boxed title style={colback=#2, arc=1.5mm, boxrule=0pt, left=2mm, right=2mm},
  coltitle=white, fonttitle=\sffamily\bfseries,
  title={#3}, #1
}
\newenvironment{definition}[1][Définition]
  {\begin{fiche}{colDef}{\ding{110}~~#1}}{\end{fiche}}
\newenvironment{propriete}[1][Propriété]
  {\begin{fiche}{colProp}{\ding{72}~~#1}}{\end{fiche}}
\newenvironment{exemple}[1][Exemple]
  {\begin{fiche}{colEx}{\ding{217}~~#1}}{\end{fiche}}
\newenvironment{remarque}[1][Remarque]
  {\begin{fiche}{colRem}{\ding{43}~~#1}}{\end{fiche}}

% Démonstration (hors boîte, style discret)
\newenvironment{demo}
  {\par\smallskip\noindent{\sffamily\bfseries\color{qtealdark}Démonstration.}\ \itshape}
  {\ \normalfont\hfill$\blacksquare$\par\medskip}

\hypersetup{colorlinks=true, linkcolor=qtealdark, urlcolor=qtealdark}

% ============================================================
\begin{document}

% ------------------- BANNIERE DE TITRE -------------------
\begin{tcolorbox}[
  enhanced, boxrule=0pt, arc=4mm,
  interior style={left color=qtealdark, right color=qteal},
  top=8mm, bottom=8mm, left=10mm, right=32mm,
  drop shadow={black!35, shadow xshift=1.5mm, shadow yshift=-1.5mm},
  before skip=0pt, after skip=9mm,
  overlay={
    \node[anchor=north west] at ([xshift=6mm,yshift=-6mm]frame.north west)
      {\includegraphics[height=1.25cm]{\logofile}};
    \node[anchor=north east, fill=qamber, text=qambertext, rounded corners=2pt,
          inner sep=4pt, font=\sffamily\bfseries\small] at ([xshift=-6mm,yshift=-6mm]frame.north east) {COURS};
  }
]
\hspace{1.6cm}\centering
{\sffamily\bfseries\color{white}\fontsize{23}{27}\selectfont \chapitretitre}\\[3mm]
{\sffamily\color{white!90}\large \matiere\ \textbullet\ \niveau}
\end{tcolorbox}

% ------------------- OBJECTIFS -------------------
\begin{tcolorbox}[
  enhanced, colback=tealpale, colframe=qteal, boxrule=0.8pt, arc=2mm,
  left=5mm, right=5mm, top=3mm, bottom=3mm,
  title={\sffamily\bfseries\color{qtealdark}\ding{51}~~Objectifs du chapitre},
  coltitle=qtealdark, colbacktitle=tealpale, boxrule=0.8pt,
  attach title to upper={\par\vspace{1mm}}
]
\begin{itemize}[leftmargin=6mm, label=\ding{212}, itemsep=1mm]
  \item Construire $\ln$ comme primitive de $x \mapsto \frac{1}{x}$ s'annulant en $1$, et en déduire ses propriétés fondamentales.
  \item Démontrer et utiliser la relation fonctionnelle et toutes les règles de calcul algébrique.
  \item Maîtriser les limites usuelles et les croissances comparées, avec démonstration des cas clés.
  \item Résoudre tous les types d'équations et d'inéquations logarithmiques, y compris avec changement de variable.
  \item Dériver des fonctions composées avec $\ln u$, étudier des fonctions complètes, utiliser le logarithme décimal.
  \item Traiter les exercices type examen national : étude de fonction avec $\ln$, suite liée à $\ln$, inégalités classiques.
\end{itemize}
\end{tcolorbox}

\vspace{4mm}

% ------------------- SECTION 1 -------------------
\section{Construction et premières propriétés}

\begin{definition}[Logarithme népérien]
La fonction \textbf{logarithme népérien}, notée $\ln$, est l'unique primitive sur $]0\,;+\infty[$ de la fonction $x \mapsto \dfrac{1}{x}$ qui s'annule en $1$ :
\[
  \ln : \ ]0\,;+\infty[\ \to \mathbb{R}, \qquad \ln(1) = 0, \qquad (\forall x > 0)\ \ln'(x) = \frac{1}{x}.
\]
(L'existence de cette primitive est garantie par la continuité de $x \mapsto \frac{1}{x}$ sur $]0\,;+\infty[$ ; l'unicité, par la propriété selon laquelle deux primitives d'une même fonction sur un intervalle diffèrent d'une constante.)
\end{definition}

\begin{propriete}[Conséquences immédiates]
\begin{itemize}[leftmargin=6mm, itemsep=0.5mm]
  \item $\ln$ est continue et dérivable sur $]0\,;+\infty[$, avec $\ln'(x) = \frac{1}{x} > 0$ : $\ln$ est \textbf{strictement croissante}.
  \item $\ln x = 0 \iff x = 1$ ; $\ \ln x > 0 \iff x > 1$ ; $\ \ln x < 0 \iff 0 < x < 1$ (conséquences de la stricte croissance et de $\ln 1 = 0$).
  \item Le réel $e$ est défini par $\ln e = 1$ ; on admet $e \approx 2{,}71828$.
  \item Pour $a, b > 0$ : $\ln a = \ln b \iff a = b$, et $\ln a < \ln b \iff a < b$ (stricte croissance = injectivité + comparaison).
\end{itemize}
\end{propriete}

\begin{remarque}[Le réflexe domaine de définition]
Avant toute équation, inéquation, ou étude de fonction avec $\ln$, détermine \textbf{en premier} l'ensemble des $x$ pour lesquels les expressions sous les logarithmes sont \textbf{strictement positives}. Toute solution hors de ce domaine doit être rejetée -- c'est l'erreur la plus fréquente au bac.
\end{remarque}

% ------------------- SECTION 2 -------------------
\section{Relation fonctionnelle et règles de calcul}

\begin{propriete}[Relation fonctionnelle fondamentale]
Pour tous $a, b > 0$ :
\[
  \ln(ab) = \ln a + \ln b.
\]
\end{propriete}

\begin{demo}
Fixons $b > 0$ et posons $g(x) = \ln(bx) - \ln x$ pour $x > 0$. Alors $g$ est dérivable et
\[
  g'(x) = \frac{b}{bx} - \frac{1}{x} = \frac{1}{x} - \frac{1}{x} = 0
\]
sur $]0\,;+\infty[$ : $g$ est donc \textbf{constante}. On calcule cette constante en $x = 1$ : $g(1) = \ln b - \ln 1 = \ln b$. Donc pour tout $x > 0$ : $\ln(bx) - \ln x = \ln b$, soit $\ln(bx) = \ln x + \ln b$. En posant $a = x$, on obtient $\ln(ab) = \ln a + \ln b$.
\end{demo}

\begin{propriete}[Règles de calcul dérivées]
Pour tous $a, b > 0$, $n \in \mathbb{Z}$, $r \in \mathbb{Q}$ :
\[
  \ln\frac{1}{b} = -\ln b, \qquad
  \ln\frac{a}{b} = \ln a - \ln b, \qquad
  \ln\left(a^n\right) = n \ln a, \qquad
  \ln\sqrt{a} = \frac{1}{2}\ln a, \qquad
  \ln\left(a^r\right) = r \ln a.
\]
\end{propriete}

\begin{demo}
$\ln\frac{1}{b}$ : comme $b \times \frac{1}{b} = 1$, la relation fonctionnelle donne $\ln b + \ln\frac{1}{b} = \ln 1 = 0$, d'où $\ln\frac{1}{b} = -\ln b$. Puis $\ln\frac{a}{b} = \ln\left(a \times \frac{1}{b}\right) = \ln a + \ln\frac{1}{b} = \ln a - \ln b$. Pour $\ln(a^n)$ ($n \in \mathbb{N}$), récurrence immédiate à partir de la relation fonctionnelle ; le cas $n$ négatif s'en déduit avec $\ln\frac{1}{b}$. Enfin $\ln\sqrt{a} = \ln\left(a^{1/2}\right)$ vérifie $2\ln\sqrt{a} = \ln\left((\sqrt{a})^2\right) = \ln a$, d'où le résultat.
\end{demo}

\begin{exemple}[Simplification d'expressions]
Simplifions $A = \ln 8 + \ln 3 - \ln 6$ et $B = \ln\bigl(\sqrt{2}+1\bigr) + \ln\bigl(\sqrt{2}-1\bigr)$.
\[
  A = \ln\frac{8\times 3}{6} = \ln 4 = \ln 2^2 = 2\ln 2, \qquad
  B = \ln\bigl[(\sqrt2+1)(\sqrt2-1)\bigr] = \ln(2-1) = \ln 1 = 0.
\]
\end{exemple}

% ------------------- SECTION 3 -------------------
\section{Étude complète et limites}

\begin{propriete}[Limites aux bornes]
\[
  \lim_{x \to +\infty} \ln x = +\infty, \qquad \lim_{x \to 0^+} \ln x = -\infty.
\]
\end{propriete}

\begin{demo}
Pour $\lim\limits_{x\to+\infty}\ln x = +\infty$ : montrons d'abord que $\ln(2^n) = n\ln 2 \to +\infty$ quand $n \to +\infty$ (car $\ln 2 > 0$, $2 > 1$). Comme $\ln$ est croissante, pour tout $x \geq 2^n$, $\ln x \geq n \ln 2$ : donc $\ln x \to +\infty$ quand $x \to +\infty$. Pour la limite en $0^+$, on pose $x = \frac1t$ avec $t \to +\infty$ : $\ln x = \ln\frac1t = -\ln t \to -\infty$.
\end{demo}

\begin{remarque}[Asymptote et branche parabolique]
La courbe de $\ln$ admet l'axe des ordonnées pour \textbf{asymptote verticale}. En $+\infty$, on verra que $\frac{\ln x}{x} \to 0$ : la courbe présente une \textbf{branche parabolique de direction $(Ox)$}.
\end{remarque}

\begin{propriete}[Croissances comparées et limites usuelles]
\[
  \lim_{x \to +\infty} \frac{\ln x}{x} = 0, \qquad
  \lim_{x \to +\infty} \frac{\ln x}{x^n} = 0 \ (n \in \mathbb{N}^*), \qquad
  \lim_{x \to 0^+} x \ln x = 0, \qquad
  \lim_{x \to 0^+} x^n \ln x = 0,
\]
\[
  \lim_{x \to 0} \frac{\ln(1+x)}{x} = 1, \qquad
  \lim_{x \to 1} \frac{\ln x}{x-1} = 1.
\]
\end{propriete}

\begin{demo}
Montrons $\lim\limits_{x\to+\infty}\frac{\ln x}{x} = 0$. Pour $t \geq 1$, on a $\ln t \leq t - 1 \leq t$ (inégalité classique, obtenue via la concavité de $\ln$, voir plus loin). En appliquant cette inégalité à $\sqrt{x}$ (pour $x \geq 1$) : $\ln\sqrt{x} \leq \sqrt{x}$, soit $\frac{1}{2}\ln x \leq \sqrt x$, donc $0 < \frac{\ln x}{x} \leq \frac{2\sqrt x}{x} = \frac{2}{\sqrt x} \to 0$. Par gendarmes, $\frac{\ln x}{x} \to 0$.

\smallskip
Pour $\lim\limits_{x\to0^+} x\ln x = 0$ : poser $t = \frac1x \to +\infty$, alors $x\ln x = \frac1t \ln\frac1t = -\frac{\ln t}{t} \to 0$ (d'après ce qui précède).

\smallskip
Pour $\lim\limits_{x\to0}\frac{\ln(1+x)}{x} = 1$ : c'est le taux d'accroissement de $\ln$ en $1$, puisque $\frac{\ln(1+x)-\ln(1)}{x} \to \ln'(1) = 1$ par définition du nombre dérivé.
\end{demo}

\begin{exemple}[Calculs de limites avec croissances comparées]
\textbf{a)} $\displaystyle\lim_{x \to +\infty} \frac{\ln x - x}{x} = \lim_{x\to+\infty}\left(\frac{\ln x}{x} - 1\right) = 0 - 1 = -1$.

\smallskip
\textbf{b)} $\displaystyle\lim_{x \to 0^+} x^2 \ln x = 0$ (croissance comparée directe).

\smallskip
\textbf{c)} $\displaystyle\lim_{x \to +\infty} \bigl(\ln(x+1) - \ln x\bigr) = \lim_{x\to+\infty}\ln\frac{x+1}{x} = \ln 1 = 0$ (par continuité de $\ln$ en $1$, car $\frac{x+1}{x}\to 1$).
\end{exemple}

\begin{propriete}[Concavité et inégalité fondamentale]
$\ln$ est \textbf{concave} sur $]0\,;+\infty[$ (car $\ln''(x) = -\frac{1}{x^2} < 0$). La tangente en $x=1$ a pour équation $y = x - 1$, et par concavité la courbe est en dessous de cette tangente :
\[
  \ln x \leq x - 1 \qquad \text{pour tout } x > 0 \quad (\text{égalité si et seulement si } x = 1).
\]
\end{propriete}

% ------------------- SECTION 4 -------------------
\section{Équations, inéquations : tous les cas}

\begin{exemple}[Équation directe]
Résolvons $\ln(x-1) + \ln(x+2) = \ln 4$.

\smallskip
\textbf{Domaine} : $x-1>0$ et $x+2>0$, soit $x>1$ : $D=\,]1\,;+\infty[$.

\smallskip
Sur $D$ : $\ln\bigl[(x-1)(x+2)\bigr]=\ln4 \iff (x-1)(x+2)=4 \iff x^2+x-6=0 \iff (x-2)(x+3)=0$. Racines $x=2$, $x=-3$ ; seule $x=2\in D$ convient.

\smallskip
\textbf{Conclusion} : $S=\{2\}$.
\end{exemple}

\begin{exemple}[Inéquation directe]
Résolvons $\ln(2x-1)\leq 0$.

\smallskip
\textbf{Domaine} : $x>\frac12$. Sur ce domaine : $\ln(2x-1)\leq \ln1 \iff 2x-1\leq1 \iff x\leq1$ (stricte croissance).

\smallskip
\textbf{Conclusion} : $S=\left]\frac12\,;1\right]$.
\end{exemple}

\begin{exemple}[Équation avec changement de variable]
Résolvons $(\ln x)^2 - 3\ln x + 2 = 0$.

\smallskip
\textbf{Domaine} : $x>0$. On pose $t=\ln x$ (sans restriction de signe, $t\in\mathbb{R}$) : $t^2-3t+2=0 \iff (t-1)(t-2)=0$, donc $t=1$ ou $t=2$.

\smallskip
Retour à $x$ : $\ln x=1 \iff x=e$ ; $\ln x=2 \iff x=e^2$.

\smallskip
\textbf{Conclusion} : $S=\{e\,;e^2\}$.
\end{exemple}

\begin{exemple}[Système avec somme et produit]
Résoudre le système $\begin{cases} \ln x + \ln y = \ln 6 \\ x + y = 5 \end{cases}$ (avec $x,y>0$).

\smallskip
La première équation équivaut à $xy = 6$ (domaine $x,y>0$ assuré). Donc $x$ et $y$ sont les racines de $X^2 - 5X + 6 = 0$, soit $(X-2)(X-3)=0$.

\smallskip
\textbf{Conclusion} : $\{x,y\}=\{2,3\}$, soit $(x,y)=(2,3)$ ou $(3,2)$.
\end{exemple}

\begin{remarque}[Piège : équation avec valeur absolue implicite]
$\ln(x^2) = 2\ln|x|$ (et \textbf{non} $2\ln x$) car $x^2>0$ dès que $x\neq0$, y compris pour $x<0$, alors que $\ln x$ n'est pas défini pour $x<0$. Toujours vérifier si l'expression sous le logarithme peut être positive pour des $x$ de signe quelconque.
\end{remarque}

% ------------------- SECTION 5 -------------------
\section{Dérivée de $\ln u$, primitives, logarithme décimal}

\begin{propriete}[Dérivée de la composée]
Si $u$ est dérivable et \textbf{strictement positive} sur $I$ :
\[
  \bigl(\ln u\bigr)' = \frac{u'}{u}.
\]
Conséquence pour les primitives : les primitives de $\dfrac{u'}{u}$ sont $\ln|u| + C$ (valable même si $u<0$, à condition que $u$ ne s'annule pas).
\end{propriete}

\begin{demo}
Composée de $u$ (dérivable en $x$) et de $\ln$ (dérivable en $u(x) > 0$, de dérivée $\frac{1}{u(x)}$) : $(\ln u)'(x) = u'(x)\times\frac{1}{u(x)} = \frac{u'(x)}{u(x)}$.
\end{demo}

\begin{definition}[Logarithme décimal]
Pour $x>0$, le \textbf{logarithme décimal} est $\log x = \dfrac{\ln x}{\ln 10}$. Il vérifie $\log 1=0$, $\log 10=1$, $\log(10^n)=n$ pour $n\in\mathbb{Z}$, et les mêmes règles de calcul que $\ln$ (produit, quotient, puissance). \emph{Usage} : mesure du pH en chimie ($\text{pH}=-\log[\text{H}_3\text{O}^+]$), échelle de Richter, etc.
\end{definition}

\begin{exemple}[Étude complète d'une fonction avec $\ln$]
Étudions $f(x) = x - \ln x$ sur $D_f=\,]0\,;+\infty[$ : variations, limites, tableau.

\smallskip
\textbf{Dérivée} : $f'(x) = 1 - \dfrac1x = \dfrac{x-1}{x}$. Signe de $x-1$ (car $x>0$) : $f'<0$ sur $]0\,;1[$, $f'>0$ sur $]1\,;+\infty[$ : \textbf{minimum} $f(1) = 1 - 0 = 1$.

\smallskip
\textbf{Limites} : en $0^+$ : $x\to0$ et $-\ln x\to+\infty$, donc $f(x)\to+\infty$. En $+\infty$ : $f(x)=x\left(1-\frac{\ln x}{x}\right)\to+\infty$ (car $\frac{\ln x}{x}\to0$).

\smallskip
\textbf{Conséquence} : comme $f(x)\geq f(1)=1>0$ pour tout $x>0$, on retrouve l'inégalité $x > \ln x$ pour tout $x>0$ (plus forte que $\ln x \leq x-1$ pour $x\ne1$, elles sont en fait cohérentes).
\end{exemple}

\begin{exemple}[Dérivée d'une composée et étude de signe]
Dérivons $f(x)=\ln\left(x^2+1\right)$ sur $\mathbb{R}$ (possible car $x^2+1>0$ toujours) :
\[
  f'(x) = \frac{2x}{x^2+1}.
\]
Signe de $f'$ = signe de $x$ : $f$ décroît sur $]-\infty\,;0]$, croît sur $[0\,;+\infty[$, minimum $f(0)=\ln1=0$.
\end{exemple}

% ------------------- SECTION 6 -------------------
\section{Exercices corrigés -- niveau examen national SM}

\begin{exemple}[Exercice 1 -- domaine, dérivée, limites]
Soit $f(x) = \ln\left(\dfrac{x-1}{x+1}\right)$. Déterminer $D_f$, calculer $f'(x)$, et les limites aux bornes.

\smallskip
\textbf{Corrigé.} \emph{Domaine} : $\dfrac{x-1}{x+1}>0 \iff x\in\,]-\infty\,;-1[\,\cup\,]1\,;+\infty[$ (tableau de signes).

\smallskip
\emph{Dérivée} : posons $u(x)=\dfrac{x-1}{x+1}$, avec $u'(x) = \dfrac{2}{(x+1)^2}$ (dérivée d'un quotient). Alors $f=\ln u$, donc $f'=\dfrac{u'}{u}$ :
\[
  f'(x) = \frac{2}{(x+1)^2}\times\frac{x+1}{x-1} = \frac{2}{(x+1)(x-1)} = \frac{2}{x^2-1}.
\]

\smallskip
\emph{Limites} : en $-\infty$ et $+\infty$ : $\frac{x-1}{x+1}\to1$, donc $f(x)\to\ln1=0$. En $1^+$ : $\frac{x-1}{x+1}\to0^+$, donc $f(x)\to-\infty$. En $(-1)^-$ : $\frac{x-1}{x+1}\to+\infty$ (num. $\to-2$, dénom. $\to0^-$), donc $f(x)\to+\infty$.
\end{exemple}

\begin{exemple}[Exercice 2 -- suite définie via $\ln$]
Soit $u_n=\ln(n^2+1)-\ln(n^2)$ pour $n\geq1$. Simplifier $u_n$ et calculer $\lim u_n$.

\smallskip
\textbf{Corrigé.} $u_n=\ln\dfrac{n^2+1}{n^2}=\ln\left(1+\dfrac{1}{n^2}\right)$. Comme $\dfrac{1}{n^2}\to0$ et $\ln$ continue en $1$ (avec $\ln1=0$) : $u_n\to\ln1=0$.

\smallskip
\emph{Variante plus fine} : en écrivant $u_n = \dfrac{\ln(1+1/n^2)}{1/n^2}\times\dfrac{1}{n^2}$, et en utilisant $\frac{\ln(1+x)}{x}\to1$ (avec $x=\frac{1}{n^2}\to0$), on obtient $u_n \sim \frac{1}{n^2}$, ce qui précise la vitesse de convergence vers $0$.
\end{exemple}

\begin{exemple}[Exercice 3 -- inéquation avec quotient]
Résoudre $\ln\left(\dfrac{1}{x}\right) \geq \ln(x-2)$.

\smallskip
\textbf{Corrigé.} \emph{Domaine} : $\frac1x>0 \iff x>0$, et $x-2>0 \iff x>2$. Donc $D=\,]2\,;+\infty[$.

\smallskip
Sur $D$, par stricte croissance de $\ln$ : $\dfrac{1}{x}\geq x-2 \iff 1 \geq x(x-2)$ (car $x>0$) $\iff x^2-2x-1\leq0$. Racines de $X^2-2X-1$ : $X=1\pm\sqrt2$. Donc $x^2-2x-1\leq0 \iff x\in[1-\sqrt2\,;1+\sqrt2]$.

\smallskip
En intersectant avec $D=\,]2\,;+\infty[$ : comme $1+\sqrt2\approx2{,}41$, on obtient $S=\,]2\,;1+\sqrt2]$.
\end{exemple}

\begin{exemple}[Exercice 4 -- démontrer une inégalité par étude de fonction]
Montrer que pour tout $x>0$ : $\ln x \leq 2(\sqrt x - 1)$.

\smallskip
\textbf{Corrigé.} Posons $g(x) = 2(\sqrt x-1) - \ln x$ sur $]0\,;+\infty[$. $g'(x) = \dfrac{1}{\sqrt x} - \dfrac1x = \dfrac{\sqrt x - 1}{x}$ (en multipliant : $\frac{1}{\sqrt x}=\frac{\sqrt x}{x}$).

\smallskip
Signe de $g'$ = signe de $\sqrt x - 1$ : $g$ décroît sur $]0\,;1]$, croît sur $[1\,;+\infty[$ : \textbf{minimum} $g(1) = 2(1-1)-\ln1 = 0$.

\smallskip
Donc $g(x)\geq g(1) = 0$ pour tout $x>0$, c'est-à-dire $\ln x \leq 2(\sqrt x-1)$.
\end{exemple}

\begin{exemple}[Exercice 5 -- limite avec composée exponentielle-log implicite]
Calculer $\displaystyle\lim_{x\to+\infty}\bigl(\ln(2x+1)-\ln(x+3)\bigr)$ puis $\displaystyle\lim_{x\to0^+} x^2\ln(3x)$.

\smallskip
\textbf{Corrigé.} \emph{Premier calcul} : $\ln(2x+1)-\ln(x+3)=\ln\dfrac{2x+1}{x+3}\to\ln2$ (car $\frac{2x+1}{x+3}\to2$ et $\ln$ continue en $2$).

\smallskip
\emph{Second calcul} : $x^2\ln(3x) = x^2\ln3 + x^2\ln x \to 0 + 0 = 0$ (croissance comparée $x^2\ln x\to0$ en $0^+$, et $x^2\ln3\to0$).
\end{exemple}

% ------------------- SECTION 7 -------------------
\section{Méthodes et pièges : la boîte à outils de l'examen}

\subsection{Ce qu'on te demande, ce que tu fais}

\begin{center}
\renewcommand{\arraystretch}{1.45}
\sffamily\small
\begin{tabular}{>{\raggedright\arraybackslash}p{7.2cm} >{\raggedright\arraybackslash}p{8cm}}
\rowcolor{qtealdark} \color{white}\bfseries Ce qu'on te demande & \color{white}\bfseries Ton réflexe \\
\rowcolor{tealpale} Équation/inéquation avec $\ln$ & Domaine d'abord (arguments $>0$), puis simplifier avec les règles de calcul. \\
Équation type $(\ln x)^2+\cdots$ & Changement de variable $t=\ln x$, résoudre en $t$, revenir à $x=e^t$. \\
\rowcolor{tealpale} Limite avec $\ln$ en $+\infty$ ou $0^+$ & Identifier une croissance comparée : $\frac{\ln x}{x^n}\to0$, $x^n\ln x\to0$. \\
Limite du type $\frac{\ln(1+f(x))}{f(x)}$ avec $f(x)\to0$ & Reconnaître $\to1$ (taux d'accroissement de $\ln$ en $1$). \\
\rowcolor{tealpale} Dériver $\ln u$ & $(\ln u)' = \frac{u'}{u}$, vérifier $u>0$ sur le domaine. \\
Primitive d'un quotient $\frac{u'}{u}$ & $\ln|u| + C$. \\
\rowcolor{tealpale} Démontrer une inégalité avec $\ln$ & Étudier le signe de (membre de droite) $-$ (membre de gauche) via sa dérivée. \\
Étude complète d'une fonction avec $\ln$ & Domaine, limites (croissances comparées si besoin), dérivée, tableau, asymptotes. \\
\end{tabular}
\end{center}

\subsection{Les pièges qui coûtent des points}

\begin{remarque}[Erreurs relevées chaque année par les correcteurs]
\begin{itemize}[leftmargin=6mm, itemsep=0.5mm]
  \item Résoudre une équation/inéquation avec $\ln$ \textbf{sans déterminer le domaine} au préalable.
  \item Écrire $\ln(x^2) = 2\ln x$ au lieu de $2\ln|x|$ (faux si $x$ peut être négatif).
  \item Écrire $\ln(a+b) = \ln a + \ln b$ : \textbf{faux}, la relation ne s'applique qu'au \textbf{produit}.
  \item Oublier de vérifier le signe de $u$ avant d'écrire $(\ln u)' = \frac{u'}{u}$.
  \item Confondre $\ln x \leq x - 1$ (vraie pour tout $x>0$) avec une égalité générale.
  \item Dans un changement de variable $t=\ln x$ : oublier que $t$ parcourt \textbf{tout} $\mathbb{R}$ (pas de restriction de signe sur $t$, contrairement à $x$).
  \item Confondre $\ln$ (logarithme népérien, base $e$) et $\log$ (logarithme décimal, base $10$) dans un exercice de chimie ou physique.
\end{itemize}
\end{remarque}

% ------------------- RESUME -------------------
\vspace{4mm}
\begin{tcolorbox}[
  enhanced, boxrule=0pt, arc=2mm,
  interior style={left color=qtealdark, right color=qteal},
  left=6mm, right=6mm, top=4mm, bottom=4mm,
  fontupper=\color{white}\sffamily,
  title={\sffamily\bfseries\color{white}\ding{72}~~À retenir},
  colbacktitle=qtealdark, coltitle=white, boxrule=0pt
]
\begin{itemize}[leftmargin=6mm, label={\color{qamber}\ding{212}}, itemsep=1mm]
  \item $\ln$ : primitive de $\frac1x$ sur $]0\,;+\infty[$, $\ln1=0$, $\ln e=1$, strictement croissante ; $\ln(ab)=\ln a+\ln b$ (démontrée).
  \item Limites : $\ln x\to+\infty$ ($x\to+\infty$), $\ln x\to-\infty$ ($x\to0^+$) ; $\frac{\ln x}{x^n}\to0$ ; $x^n\ln x\to0$ ; $\frac{\ln(1+x)}{x}\to1$.
  \item Concavité : $\ln x \leq x-1$ pour tout $x>0$.
  \item $(\ln u)'=\frac{u'}{u}$ ; primitives de $\frac{u'}{u}$ : $\ln|u|+C$ ; $\log x=\frac{\ln x}{\ln10}$.
  \item Toujours commencer par le \textbf{domaine} ; changement de variable $t=\ln x$ pour les équations polynomiales en $\ln x$.
\end{itemize}
\end{tcolorbox}

\end{document}
